\documentclass[11pt]{article}
\usepackage[margin=1in]{geometry}
\usepackage{booktabs}
\usepackage{amsmath}
\usepackage{url}
\usepackage[hidelinks]{hyperref}

\title{A Census of the Model Context Protocol Registry}

\author{Ashish Sinha\\
\texttt{sinha.ashish.4.u@gmail.com}}

\date{}

\begin{document}
\maketitle

\begin{abstract}
The Model Context Protocol (MCP) has acquired a public registry of servers that
agents can discover and connect to, but the contents of that registry have not
been described. We retrieve it in full---$82{,}994$ version records resolving
to $25{,}125$ distinct servers from $15{,}468$ publishers---and report what is
in it.

Three findings. First, the registry is a version log rather than a catalogue,
and the distinction changes conclusions by an order of magnitude: $74.5\%$ of
version records share a description with another record, but only $3.6\%$ of
distinct servers do. Treating version history as duplication overstates it by
$20.7\times$. One server accounts for $1{,}175$ published versions.

Second, $393$ servers ($1.56\%$) declare neither a remote endpoint nor an
installable package. They are listed and discoverable but cannot be connected
to or installed by any means the registry describes.

Third, the publisher distribution is extremely long-tailed: $89.5\%$ of
publishers have registered exactly one server, while the ten largest account
for $15.7\%$ of all servers. Registry-level statistics are therefore
dominated by a handful of bulk publishers, and any measure that does not
account for this describes those publishers rather than the ecosystem.

We also find $5{,}643$ servers ($22.5\%$) with no repository link and eight
still shipping the unedited template string \texttt{[describe what your server
does]}. Retrieval and analysis code is released.
\end{abstract}

\section{Introduction}

The Model Context Protocol standardises how language-model agents connect to
external tools. A public registry lists servers implementing it, and agent
runtimes increasingly consult that registry to decide what an agent can reach.
What the registry actually contains has not, to our knowledge, been measured.

This matters for a practical reason. When a registry is used for discovery, its
contents become the search space an agent selects from. Entries that cannot be
installed, that are indistinguishable from one another, or that are duplicated
by version churn all degrade that selection, and none of them are visible from
a single listing page.

We retrieve the registry in full and describe it. Our contributions are a
complete census as of August 2026, a methodological caution about the unit of
analysis that changes headline numbers by $20.7\times$, and released code that
reproduces every figure.

\section{Data}

We retrieved every record from the public registry endpoint
\texttt{registry.modelcontextprotocol.io/v0/servers} by cursor pagination in
August 2026, continuing until the endpoint returned no further cursor:
$82{,}994$ records over $830$ pages.

Each record carries a server name, a description, a version, a publication
status, and an \texttt{isLatest} flag, together with optional
\texttt{repository}, \texttt{remotes}, \texttt{packages}, \texttt{title},
\texttt{websiteUrl}, and \texttt{icons} fields. Names are namespaced, so the
segment before the first slash identifies the publisher.

\subsection{The unit of analysis}

The endpoint returns one record per published \emph{version}, not per server.
Filtering to \texttt{isLatest} leaves $25{,}125$ distinct servers, a mean of
$3.30$ versions each. The maximum is $1{,}175$ versions of a single server.

This distinction is not cosmetic. Table~\ref{tab:unit} shows the same
duplicate-description measurement computed both ways.

\begin{table}[h]
\centering
\small
\begin{tabular}{lrr}
\toprule
Unit of analysis & Records & Sharing a description \\
\midrule
Version records (as returned) & 82{,}994 & 61{,}859 \quad (74.5\%) \\
Distinct servers (\texttt{isLatest}) & 25{,}125 & 904 \quad (3.6\%) \\
\midrule
\multicolumn{2}{l}{Overstatement factor} & $\mathbf{20.7\times}$ \\
\bottomrule
\end{tabular}
\caption{Duplicate descriptions under two units of analysis. A version log read
as a catalogue reports a registry that is three-quarters duplicated; it is not.}
\label{tab:unit}
\end{table}

Read naively, the registry appears to be $74.5\%$ duplicate. It is not. Nearly
all of that figure is one publisher republishing one server, and a measurement
that does not deduplicate to latest versions is measuring release cadence. All
figures below use distinct servers.

\section{Findings}

\subsection{Servers that cannot be reached}

A registry entry is actionable only if it says how to connect. A server
declares either \texttt{remotes} (a hosted endpoint) or \texttt{packages} (an
installable artefact), or both.

\begin{table}[h]
\centering
\small
\begin{tabular}{lrr}
\toprule
Transport & Servers & Share \\
\midrule
Remote, streamable HTTP & 13{,}279 & 52.9\% \\
Package, npm            & 7{,}985  & 31.8\% \\
Package, PyPI           & 3{,}434  & 13.7\% \\
Remote, SSE             & 1{,}054  & 4.2\%  \\
Package, MCPB           & 891      & 3.5\%  \\
Package, OCI            & 821      & 3.3\%  \\
Package, NuGet          & 105      & 0.4\%  \\
\midrule
\textbf{Neither}        & \textbf{393} & \textbf{1.56\%} \\
\bottomrule
\end{tabular}
\caption{Declared transports. Servers may declare more than one, so shares
exceed $100\%$. The final row counts servers declaring none.}
\label{tab:transport}
\end{table}

$393$ servers declare no remote endpoint and no package. They occupy registry
namespace and appear in discovery results, but nothing in their record
describes how a client would reach them. A further $262$ servers are marked
deprecated but remain listed.

\subsection{Publisher concentration}

$25{,}125$ servers come from $15{,}468$ publishers, but the distribution is
severely skewed. $13{,}848$ publishers ($89.5\%$) have registered exactly one
server. The largest single publisher accounts for $1{,}312$; the ten largest
together account for $15.7\%$ of the registry.

The consequence is that unweighted registry statistics describe bulk
publishers rather than typical practice. Any claim of the form ``$X\%$ of MCP
servers do $Y$'' should be checked against whether $X$ survives removing the
top ten publishers.

\subsection{Metadata completeness}

No required field is missing from any record: name, description, version, and
status are present on all $25{,}125$ servers. Optional fields vary widely, and
we note two whose absence is consequential for anyone evaluating a server
before connecting to it.

$5{,}643$ servers ($22.5\%$) provide no \texttt{repository} link, so their
implementation cannot be inspected from the registry entry. Eight servers still
carry the unedited template string \texttt{[describe what your server does]} as
their description---a small number, but one that indicates entries reach
publication without review.

Descriptions are capped at $100$ characters, with a median length of $89$. The
cap is a registry constraint rather than an authoring choice, and it bounds how
much any consumer can learn about a server before connecting.

\subsection{Description similarity}

Applying near-duplicate detection to descriptions of distinct servers
($3$-word shingles, Jaccard, MinHash-LSH above $3{,}000$ records with exact
verification of every candidate) yields $487$ clusters at threshold $0.70$,
containing $1{,}252$ servers ($4.98\%$).

Every one of these clusters is a \emph{shared-stem variant} rather than a set
of identical records: the members share description text but differ in server
name. None are identical across all evaluated fields. Descriptions alone are
therefore a weak discriminator between registry entries---which matters
directly, because a registry consumer choosing between servers sees little
else.

Sweeping the threshold gives a notably flat curve: $19.4$ clusters per $1{,}000$
servers at $0.7$, rising only to $24.9$ at $0.4$, a $1.05\times$ increase over
the final step. Short, capped marketing copy produces little near-duplicate
structure at any threshold.

\section{Limitations}

The census is a snapshot taken in August 2026 of a registry under active
growth; the absolute counts will drift, though the structural observations
should not.

We measure only what records declare. We did not attempt to connect to any
endpoint, so we cannot say how many of the $24{,}732$ servers that \emph{do}
declare a transport are actually reachable. The $393$ figure is a lower bound
on unusable entries, established from the records alone.

MinHash-LSH offers no recall guarantee, so cluster counts are lower bounds;
every reported pair was verified with its exact Jaccard score.

We do not evaluate server quality, safety, or behaviour. This is a census of
what is listed, not an assessment of what those listings do.

\section{Conclusion}

The MCP registry contains $25{,}125$ distinct servers from $15{,}468$
publishers, the overwhelming majority of whom have published exactly one. It
is mostly complete and mostly non-redundant once read at the right unit of
analysis, which is the caveat that matters most: the endpoint returns version
history, and reading it as a catalogue turns a $3.6\%$ duplication rate into a
$74.5\%$ one.

The concrete defect worth fixing is small and specific. $393$ servers are
listed with no way to reach them, $262$ deprecated servers remain in the
listing, and eight entries were published without their template being filled
in. These are validation gaps rather than design flaws, and a registry-side
check at publication time would close all three.

\section*{Reproducibility}

Retrieval and analysis code, and the audit tooling used for the similarity
measurements, are at
\url{https://github.com/ashishsinha1602/dataset-integrity-audit}.

\begin{verbatim}
python fetch_mcp_registry.py --out mcp-registry.jsonl
python audit.py prep --data mcp-latest.jsonl --text-field description \
                     --answer-field name --group-field _status
python audit.py curve --data mcp-latest.jsonl --text-field description \
                      --answer-field name
\end{verbatim}

\end{document}
